\documentclass[11pt,a4paper]{article}
\usepackage[utf8]{inputenc}
\usepackage[ngerman]{babel}
\usepackage{amsmath,amssymb,amsthm}
\usepackage{graphicx}
\usepackage{hyperref}
\usepackage{braket}
\usepackage{booktabs}
\usepackage{geometry}
\geometry{margin=2.5cm}

\newtheorem{theorem}{Satz}
\newtheorem{lemma}[theorem]{Lemma}
\newtheorem{definition}[theorem]{Definition}
\newtheorem{proposition}[theorem]{Proposition}
\newtheorem{remark}[theorem]{Bemerkung}
\newtheorem{example}[theorem]{Beispiel}

\title{Quantenfehlerkorrektur: \\
Syndrommessung und diskrete Fehlerdiagnose}

\author{
Technische Dokumentation \\
\texttt{smooth-numbers project} \\
\texttt{Juni 2026}
}

\date{\today}

\begin{document}

\maketitle

\begin{abstract}
Diese Notiz präzisiert die mathematische Struktur der Quantenfehlerkorrektur am Beispiel des zyklischen $[\![5,1,3]\!]$-Stabilisatorcodes. Wir korrigieren zwei häufige Missverständnisse: (1) Der Code-Projektor $\mathcal{P}_{\text{Code}}$ allein korrigiert keine Fehler, sondern projiziert nur auf den Code-Unterraum; Fehlerkorrektur erfordert zusätzlich Syndrommessung und Recovery-Operatoren. (2) Die 24 Hurwitz-Einheiten bilden die binäre Tetraedergruppe $2T$, nicht die binäre Ikosaedergruppe $2I \cong \text{SL}(2,5)$, welche 120 Elemente besitzt.

Wir entwickeln ein vollständiges Rechenbeispiel für einen Phasenfehler $Z_1$ und zeigen die exakte Syndromberechnung, Fehlerdiagnose und Recovery-Operation.
\end{abstract}

\section{Einleitung}

\subsection{Motivation}

Die Quantenfehlerkorrektur ist ein zentrales Konzept für die Realisierung skalierbarer Quantencomputer. Ein verbreitetes Missverständnis lautet:

\begin{quote}
\textit{„Der Code-Projektor $\mathcal{P}_{\text{Code}}$ filtert Rauschen heraus und stellt den fehlerfreien Zustand wieder her."}
\end{quote}

Diese Vorstellung ist \textbf{unvollständig}. Der Projektor allein kann keine Fehler korrigieren, da er lediglich auf den Code-Unterraum projiziert. Die eigentliche Fehlerkorrektur erfordert drei Schritte:

\begin{enumerate}
\item \textbf{Syndrommessung:} Bestimmung der Eigenwerte der Stabilisatorgeneratoren
\item \textbf{Fehlerdiagnose:} Zuordnung des Syndroms zu einem diskreten Fehler
\item \textbf{Recovery:} Anwendung des inversen Fehleroperators
\end{enumerate}

Diese Notiz entwickelt diese Struktur explizit am $[\![5,1,3]\!]$-Code.

\subsection{Algebraische Präzisierungen}

\subsubsection{Hurwitz-Einheiten und binäre Tetraedergruppe}

Die \textbf{Hurwitz-Quaternionen} sind definiert als
\[
\mathcal{H} = \left\{ a + bi + cj + dk \mid a,b,c,d \in \mathbb{Z} \text{ oder } a,b,c,d \in \mathbb{Z} + \tfrac{1}{2}, \; a+b+c+d \in 2\mathbb{Z} \right\}.
\]

Die \textbf{Hurwitz-Einheiten} (Elemente mit Norm 1) bilden eine endliche Gruppe der Ordnung 24:
\[
\{\pm 1, \pm i, \pm j, \pm k, \tfrac{1}{2}(\pm 1 \pm i \pm j \pm k)\}.
\]

Diese Gruppe ist isomorph zur \textbf{binären Tetraedergruppe} $2T$, der doppelten Überlagerung der Tetraedergruppe $T \cong A_4$:
\[
2T \cong \text{SL}(2,3), \quad |2T| = 24.
\]

\textbf{Nicht zu verwechseln} mit der \textbf{binären Ikosaedergruppe}:
\[
2I \cong \text{SL}(2,5), \quad |2I| = 120.
\]

Die binäre Ikosaedergruppe ist die doppelte Überlagerung der Ikosaedergruppe $I \cong A_5$ und hat 120 Elemente, nicht 24.

\section{Der zyklische $[\![5,1,3]\!]$-Stabilisatorcode}

\subsection{Stabilisatorgeneratoren}

Der zyklische $[\![5,1,3]\!]$-Code kodiert 1 logisches Qubit in 5 physischen Qubits und korrigiert 1 beliebigen Fehler. Die vier Stabilisatorgeneratoren lauten:

\begin{align}
S_1 &= X_1 Z_2 Z_3 X_4 I_5 = XZZXI, \\
S_2 &= I_1 X_2 Z_3 Z_4 X_5 = IXZZX, \\
S_3 &= X_1 I_2 X_3 Z_4 Z_5 = XIXZZ, \\
S_4 &= Z_1 X_2 I_3 X_4 Z_5 = ZXIXZ.
\end{align}

\subsection{Code-Unterraum}

Der \textbf{Code-Unterraum} $\mathcal{C}$ ist der gemeinsame $+1$-Eigenraum aller vier Generatoren:
\[
\mathcal{C} = \{\ket{\psi} \in (\mathbb{C}^2)^{\otimes 5} \mid S_m \ket{\psi} = \ket{\psi} \text{ für alle } m = 1,2,3,4\}.
\]

Der Unterraum ist 2-dimensional (1 logisches Qubit) und wird durch die logischen Basiszustände $\ket{0_L}, \ket{1_L}$ aufgespannt.

\subsection{Code-Projektor}

Der \textbf{Code-Projektor} ist definiert als
\[
\mathcal{P}_{\text{Code}} = \frac{1}{16} (I + S_1)(I + S_2)(I + S_3)(I + S_4).
\]

\begin{remark}[Wichtig]
$\mathcal{P}_{\text{Code}}$ projiziert auf den Code-Unterraum $\mathcal{C}$, korrigiert aber \textbf{keine Fehler}. Ein fehlerhafter Zustand wird einfach auf den Code-Unterraum projiziert, was die Fehlerinformation zerstört, aber nicht den ursprünglichen logischen Zustand wiederherstellt.
\end{remark}

\section{Fehlerkorrektur: Phasenfehler $Z_1$}

\subsection{Ausgangszustand und Fehler}

Sei $\ket{\psi} \in \mathcal{C}$ ein gültiger Codezustand mit
\[
S_m \ket{\psi} = \ket{\psi} \quad \text{für alle } m = 1,2,3,4.
\]

Wir betrachten einen \textbf{Phasenfehler} auf dem ersten Qubit:
\[
E = Z_1 = Z \otimes I \otimes I \otimes I \otimes I = ZIIII.
\]

Der fehlerhafte Zustand lautet
\[
\ket{\psi'} = Z_1 \ket{\psi}.
\]

\subsection{Syndrommessung}

Das \textbf{Syndrom} wird durch Messung der Stabilisatorgeneratoren bestimmt. Da $Z_1$ mit $X$ antikommutiert, aber mit $Z$ und $I$ kommutiert, erhalten wir:

\begin{align}
S_1 Z_1 \ket{\psi} &= (X_1 Z_2 Z_3 X_4 I_5)(Z_1) \ket{\psi} = -Z_1 X_1 Z_2 Z_3 X_4 I_5 \ket{\psi} = -Z_1 S_1 \ket{\psi} = -Z_1 \ket{\psi}, \\
S_2 Z_1 \ket{\psi} &= (I_1 X_2 Z_3 Z_4 X_5)(Z_1) \ket{\psi} = +Z_1 I_1 X_2 Z_3 Z_4 X_5 \ket{\psi} = +Z_1 S_2 \ket{\psi} = +Z_1 \ket{\psi}, \\
S_3 Z_1 \ket{\psi} &= (X_1 I_2 X_3 Z_4 Z_5)(Z_1) \ket{\psi} = -Z_1 X_1 I_2 X_3 Z_4 Z_5 \ket{\psi} = -Z_1 S_3 \ket{\psi} = -Z_1 \ket{\psi}, \\
S_4 Z_1 \ket{\psi} &= (Z_1 X_2 I_3 X_4 Z_5)(Z_1) \ket{\psi} = +Z_1^2 X_2 I_3 X_4 Z_5 \ket{\psi} = +Z_1 S_4 \ket{\psi} = +Z_1 \ket{\psi}.
\end{align}

Das \textbf{Syndrom} ist daher
\[
\sigma(Z_1) = (-1, +1, -1, +1).
\]

In binärer Notation (mit $-1 \mapsto 1$ und $+1 \mapsto 0$):
\[
\sigma_{\text{bin}}(Z_1) = (1, 0, 1, 0) = 1010.
\]

\subsection{Syndromprojektor}

Der \textbf{Syndromprojektor} für das Syndrom $(1,0,1,0)$ lautet
\[
\mathcal{P}_{1010} = \frac{1}{16} (I - S_1)(I + S_2)(I - S_3)(I + S_4).
\]

Dieser Projektor extrahiert den Unterraum aller Zustände mit diesem spezifischen Syndrom.

\subsection{Recovery-Operation}

Die Fehlerkorrektur besteht aus zwei Schritten:

\begin{enumerate}
\item \textbf{Syndrommessung:} Der Zustand $\ket{\psi'} = Z_1 \ket{\psi}$ wird durch die Messung in den Fehlerunterraum mit Syndrom $1010$ projiziert.
\item \textbf{Recovery:} Der Recovery-Operator $R = Z_1$ wird angewendet:
\[
R \ket{\psi'} = Z_1 (Z_1 \ket{\psi}) = Z_1^2 \ket{\psi} = \ket{\psi}.
\]
\end{enumerate}

Da $Z_1^2 = I$, wird der ursprüngliche Zustand exakt wiederhergestellt.

\subsection{Vollständige Fehlerkorrektur}

Die vollständige Fehlerkorrekturprozedur lautet:

\begin{equation}
\boxed{
\begin{aligned}
\ket{\psi} \quad &\xrightarrow{\text{Fehler } E = Z_1} \quad Z_1 \ket{\psi} \\
&\xrightarrow{\text{Syndrommessung}} \quad \text{Syndrom } 1010 \\
&\xrightarrow{\text{Fehlerdiagnose}} \quad \text{Fehler } Z_1 \text{ identifiziert} \\
&\xrightarrow{\text{Recovery } R = Z_1} \quad Z_1 (Z_1 \ket{\psi}) = \ket{\psi}.
\end{aligned}
}
\end{equation}

\section{DMN/ECN-Metapher: Diskretisierung kontinuierlicher Störungen}

\subsection{Von kontinuierlich zu diskret}

Die Quantenfehlerkorrektur leistet eine fundamentale Transformation:

\begin{center}
\textbf{Kontinuierliche physikalische Störung} $\quad \longrightarrow \quad$ \textbf{Diskretes Fehlersyndrom}
\end{center}

In der Praxis tritt nicht ein exakter Pauli-Fehler $Z_1$ auf, sondern eine kontinuierliche Phasenstörung:
\[
\rho \mapsto \rho' = (1-\epsilon)\rho + \epsilon Z_1 \rho Z_1^\dagger + \text{(andere Terme)}.
\]

Die Syndrommessung \textbf{diskretisiert} diese kontinuierliche Störung in ein endliches Fehleralphabet:
\[
\{\text{kein Fehler}, X_1, Y_1, Z_1, X_2, \ldots\}.
\]

\subsection{Der Übergang: Nicht Löschung, sondern Diagnose}

Die DMN/ECN-Metapher wird oft als „Hyperfokus löscht Rauschen aus" beschrieben. Eine präzisere Formulierung lautet:

\begin{center}
\begin{tabular}{ll}
\textbf{Falsch:} & Rauschen $\mapsto$ 0 \\
\textbf{Korrekt:} & Kontinuierliche Störung $\mapsto$ diskretes Syndrom $\mapsto$ korrigierbarer Fehler $\mapsto$ Zustand
\end{tabular}
\end{center}

Das ECN (Error Correction Network) wirkt nicht als bloßer \textit{Filter}, sondern als \textit{Syndromapparat}:

\begin{enumerate}
\item \textbf{Erkennung:} Das Syndrom identifiziert die Art der Abweichung
\item \textbf{Klassifikation:} Die Störung wird einem diskreten Fehlerraum zugeordnet
\item \textbf{Korrektur:} Die passende Rückoperation wird ausgeführt
\end{enumerate}

\subsection{Quaternionische Lesart}

Die Pauli-Operatoren $\{I, X, Y, Z\}$ entsprechen den Quaternionseinheiten $\{1, i, j, k\}$:

\begin{align}
I &\leftrightarrow 1, \\
X &\leftrightarrow i, \\
Y &\leftrightarrow -k, \\
Z &\leftrightarrow j.
\end{align}

Die Fehlerkorrektur überführt eine \textbf{kontinuierliche Phasenabweichung} (im kontinuierlichen $\text{SU}(2)$) in ein \textbf{endliches, klassifizierbares Fehleralphabet} (in der diskreten Pauli-Gruppe).

Dies ist analog zur Projektion einer reellen Zahl auf das nächste Element eines diskreten Gitters -- mit dem entscheidenden Unterschied, dass die Fehlerkorrektur \textit{ohne Informationsverlust} über den logischen Zustand erfolgt.

\section{Zusammenfassung}

\subsection{Präzisierte Aussagen}

\begin{enumerate}
\item \textbf{Code-Projektor:} $\mathcal{P}_{\text{Code}}$ allein korrigiert keine Fehler. Er projiziert nur auf den Code-Unterraum.
\item \textbf{Fehlerkorrektur:} Erfordert Syndrommessung + Fehlerdiagnose + Recovery-Operator.
\item \textbf{Hurwitz-Einheiten:} Die 24 Einheiten bilden die binäre Tetraedergruppe $2T \cong \text{SL}(2,3)$.
\item \textbf{Binäre Ikosaedergruppe:} $2I \cong \text{SL}(2,5)$ hat 120 Elemente, nicht 24.
\end{enumerate}

\subsection{Die korrekte Fehlerkorrekturkette}

\begin{equation}
\boxed{
\text{Kontinuierliche Störung} \quad \xrightarrow{\text{Messung}} \quad \text{Diskretes Syndrom} \quad \xrightarrow{\text{Diagnose}} \quad \text{Korrigierbarer Fehler} \quad \xrightarrow{\text{Recovery}} \quad \text{Rekonstruierter Zustand}
}
\end{equation}

Das ist die saubere mathematische Struktur der Quantenfehlerkorrektur.

\section*{Danksagung}

Diese Notiz entstand aus der Notwendigkeit, zwei häufige Missverständnisse in der Darstellung der Quantenfehlerkorrektur zu präzisieren. Der Fokus liegt auf mathematischer Klarheit und der expliziten Trennung von Code-Projektion und vollständiger Fehlerkorrektur.

\begin{thebibliography}{99}

\bibitem{nielsen2000}
M. A. Nielsen and I. L. Chuang, \textit{Quantum Computation and Quantum Information}, Cambridge University Press, 2000.

\bibitem{gottesman1997}
D. Gottesman, \textit{Stabilizer Codes and Quantum Error Correction}, Caltech PhD Thesis, 1997. arXiv:quant-ph/9705052.

\bibitem{laflamme1996}
R. Laflamme, C. Miquel, J. P. Paz, and W. H. Zurek, \textit{Perfect Quantum Error Correcting Code}, Physical Review Letters \textbf{77} (1996), 198--201.

\bibitem{conway1988}
J. H. Conway and D. A. Smith, \textit{On Quaternions and Octonions}, A K Peters, 2003.

\bibitem{knill1997}
E. Knill and R. Laflamme, \textit{Theory of quantum error-correcting codes}, Physical Review A \textbf{55} (1997), 900--911.

\end{thebibliography}

\end{document}
